% 原文 pp.35–39（印刷 pp.1153–1157）

% ===== 原文 p.35 =====
% [承接 p.34 末句「The considerations that follow generalize above 4 dimensions but become trivial below dimension 4).」]
维数 4）。
对庞加莱代数（Poincaré algebra）的表示、以及各种自旋的自由场作一般讨论，似乎提示红外自由（infrared-free）的理论可以有任意自旋的无质量粒子。
但在实际上，在我所熟悉的一切有趣例子中，都可以先验地（\textit{a priori}）论证：任何无质量粒子都只能有自旋 $0$、$1/2$ 或 $1$。

我们关心的大多数理论，不仅可以表述在平坦的 $\mathbb{R}^4$ 上，而且可以表述在带有一般度规 $g$ 的更一般的弯曲四维流形 $M^4$ 上。
事实上，任何构成对自然之描述的一部分的理论都具有这一性质，因为广义相对论是自然的一部分，而在广义相对论中时空是弯曲的！
在量子场论中，能够在弯曲时空上工作，意味着存在一个非常特殊的算符，称为应力张量（stress tensor）或能量—动量张量（energy-momentum tensor）$T_{\mu\nu}(x)$。
它度量理论对度规张量 $g$ 之改变的响应。
我们假定理论由拉格朗日量 $\mathcal{L}(\phi_i;g)$ 在一般的 $g$ 下表述，而该量在同时作用于 $\phi_i$ 与 $g$ 的微分同胚（diffeomorphism）下不变。
$g$ 不是理论的场之一——它是任意的，但在研究 $\phi_i$ 的经典或量子动力学时被固定——而这种微分同胚不变性意味着：理论若表述于时空 $(M,g)$ 之中，则它实际上只依赖于 $g$ 的微分同胚类。
在此框架下，应力张量定义为
\[
  T_{\mu\nu}=\frac{\delta\mathcal{L}}{\delta g^{\mu\nu}}.
\]
这显然（obviously）蕴含 $T$ 是对称张量
\[
  T^{\mu\nu}=T^{\nu\mu}.
\]

% [Tutor Note: 此处定义 $T_{\mu\nu}=\delta\mathcal{L}/\delta g^{\mu\nu}$ 省略了两处约定，照抄极易出错。(i) 常规的广义相对论约定是 $\delta S=\frac12\int d^4x\sqrt{g}\,T^{\mu\nu}\,\delta g_{\mu\nu}$，其中带 $\frac12$ 与 $\sqrt{g}$；作者的 $\mathcal{L}$ 是已含 $\sqrt{g}$ 的拉格朗日密度，且把 $\frac12$ 吸收进定义，故各式只差整体常数，与后文"比例常数为 $1$"的归一化一致。(ii) 变分变量是 $g^{\mu\nu}$ 而非 $g_{\mu\nu}$，二者由 $\delta g^{\mu\nu}=-g^{\mu\alpha}g^{\nu\beta}\delta g_{\alpha\beta}$ 相联系，换成对 $g_{\mu\nu}$ 变分会整体多一个负号；这也解释了为何左边写 $T_{\mu\nu}$（下指标）而右边对 $g^{\mu\nu}$（上指标）变分——指标位置并不对应，需用度规升降才能比较。]

$T$ 还可以证明满足
\[
  D_\mu T^{\mu\nu}=0
\]
这是微分同胚不变性的结果。
若我们的理论实际上是共形不变的（conformally invariant），则 $T$ 无迹（traceless），即 $g^{\mu\nu}T_{\mu\nu}=0$。

% [Tutor Note: 「由微分同胚不变性可得 $D_\mu T^{\mu\nu}=0$」跳过了两步：其一，变分 $\delta g_{\mu\nu}=D_\mu\xi_\nu+D_\nu\xi_\mu$ 后要分部积分，并用 $\phi_i$ 的场方程（即只在场方程成立时守恒，属"在壳"陈述）；其二，若量子测度不是微分同胚不变的，则会出现引力反常／迹反常，此时 $g^{\mu\nu}T_{\mu\nu}$ 正比于曲率不变量（其系数即 $a$-反常系数），无迹性只在经典层面或共形场论的不动点处成立。作者在后文讨论红外极限时始终假定反常可以忽略。]

拥有这样的应力张量会带来强有力的陈述，即使我们只限于 $M$ 为平坦欧氏空间的情形。
例如，去年秋天我们公理化量子场论时，曾要求庞加莱不变性，即对时空中每个 Killing 向量场 $K$ 都有守恒荷 $Q(K)$。
一个守恒的对称应力张量之存在是一个局部陈述，而它在整体上导致庞加莱不变性。
给定任意 Killing 向量场 $K$，人们利用 Killing 向量方程（即 $D_\mu K_\nu+D_\nu K_\mu=0$）再加上 $T$ 的对称性与守恒性，即可证明流
\[
  J^\nu(K)=K_\mu T^{\mu\nu}
\]
是守恒的。
于是得到守恒荷
\[
  Q(K)=\int_{\Sigma^{n-1}} *J(K)
\]
其中积分取在一个初值超曲面（例如时刻零）上。

% [Tutor Note: 由 $D_\mu T^{\mu\nu}=0$ 推出 $D_\nu J^\nu=0$ 时，"对称性"这一条不可省：$D_\nu(K_\mu T^{\mu\nu})=K_\mu D_\nu T^{\mu\nu}+(D_\nu K_\mu)T^{\mu\nu}$，第一项因守恒为零，第二项中 $D_\nu K_\mu$ 对 $\mu,\nu$ 对称，故只有 $T$ 的对称部分存活，再用 Killing 方程 $D_\nu K_\mu+D_\mu K_\nu=0$ 才得零。若 $T$ 不对称（例如把自旋流的贡献漏掉），则即使 $T$ 守恒也得不到 $J^\nu(K)$ 守恒，从而得不到正确的庞加莱荷。另外 $*J$ 是 $(n-1)$ 形式，才能在 $n-1$ 维超曲面上积分而不需要额外的体积元。]

我以上的讲述方式，仿佛必须先有拉格朗日量，才能推出应力张量 $T=\delta\mathcal{L}/\delta g$ 的存在。
尽管这是一条强有力的途径，人们也可以作更抽象的论证。
考虑任何能够对 $M$ 上任意度规表述的理论。
为了定义算符 $T(y)$，我们必须对任意指定的

% ===== 原文 p.36 =====
$M^4$ 上彼此不同、且与 $y$ 也不同的 $n$ 个点 $x_1,\dots,x_n$，以及算符 $\mathcal{O}_1,\dots,\mathcal{O}_n$，定义关联函数 $\langle\mathcal{O}_1(x_1)\cdots\mathcal{O}_n(x_n)T(y)\rangle_g$（这里的下标用以强调该关联函数依赖于度规 $g$）。
我们把它定义为 $\langle\mathcal{O}_1(x_1)\cdots\mathcal{O}_n(x_n)\rangle_g$ 对 $g$ 的导数：
\[
  \frac{\delta}{\delta g(y)}\langle\mathcal{O}_1(x_1)\cdots\mathcal{O}_n(x_n)\rangle_g
  =\langle\mathcal{O}_1(x_1)\cdots\mathcal{O}_n(x_n)T(y)\rangle_g.
\]

读者应当验证：在拉格朗日量存在的情形下，这个定义与先前的定义一致。
如此定义的 $T$ 显然是对称的；若理论只依赖于度规 $g$ 的微分同胚类，则它是守恒的。
局部量子场算符的许多性质都可以由 $T$ 的这个定义直接得到，而且有理由相信一般地它们全都是如此。

% [Tutor Note: 定义中要求 $x_1,\dots,x_n$ 与 $y$ 互不相同，并非无关紧要的细节：$\delta/\delta g(y)$ 是泛函导数，当 $y$ 与某个 $x_i$ 重合时会出现接触项（contact term），其系数依赖于算符的标度维数与重整化方案，故"$T(y)$ 的插入"与"对 $g$ 求导"只在点互异时才严格等价。忽略这一条件会凭空得到 $\delta^{(4)}(x_i-y)$ 型的额外贡献。]

无论如何，\footnote{研究量子引力时除外；量子引力的风格与固定时空中的量子场论非常不同，而后者才是本学期我们这些课程的主题。}人们在实践中通常考虑具有局部、守恒、对称（当然还有规范不变）的应力张量的理论。
正如我实质上已经指出的，任何出现在自然中的理论都具有这一性质，因为 $T$ 直接出现在 Einstein 方程中！

$T$ 的存在\footnote{S. Weinberg and E. Witten, \textit{Limits on massless particles}, Phys. Lett. B \textbf{96} (1980), 59--62.}对可能的无质量粒子给出严格的限制：我们将会看到，在具有应力张量的理论中，无质量粒子可能具有的自旋是 $0$、$1/2$ 与 $1$。
另一个类似的限制如下。
设 $J^\mu$ 是与某个“整体对称性”（global symmetry）相联系的守恒流。
于是 $J$ 在庞加莱变换下按向量变换，而守恒荷
\[
  Q(J)=\int *J
\]
是庞加莱不变的。
那么我们将看到，$Q$ 湮灭任何自旋为 $1$ 的无质量粒子。

这些证明彼此极为相似，所以我们把两种情形一并处理。
在四维中，用 $|p,j\rangle$ 记动量为 $p$、自旋为 $j$ 的无质量单粒子态（一般地 $j\in\mathbb{Z}/2$）。
又设 $|p',j\rangle$ 是同一庞加莱表示中动量不同的态。
假定 $|p,j\rangle$ 是 $Q$ 的本征态、本征值为某个 $q$，这不失任何实质性的普遍性。
考虑矩阵元 $\langle p',j|J^\nu|j,p\rangle$ 与 $\langle p',j|T^{\mu\nu}|j,p\rangle$。
后者绝不可能为零，而前者除非 $q=0$ 也不能为零。
因为在 $p'\to p$ 的极限下，由洛伦兹不变性我们有
\begin{enumerate}
  \item $\langle p',j|J^\nu|j,p\rangle\sim p^\nu$，以及
  \item $\langle p',j|T^{\mu\nu}|j,p\rangle\sim p^\mu p^\nu$。
\end{enumerate}
比例常数在第一种情形为 $q$，在第二种情形为 $1$；因为（$Q$ 与动量算符都由 $J$ 或 $T$ 积分而得）初态与末态相同的 $J$ 或 $T$ 的矩阵元度量的正是该态的电荷或动量。

% [Tutor Note: 这里"比例常数为 $1$"依赖于单粒子态的归一化约定。若取 $\langle p'|p\rangle=2E(2\pi)^3\delta^3(\mathbf{p}'-\mathbf{p})$，则 $\langle p|T^{\mu\nu}|p\rangle=p^\mu p^\nu$ 而 $\langle p|J^\nu|p\rangle=q\,p^\nu$；若改用协变归一化或除以 $2E$，常数会变，但"$T$ 的系数非零、$J$ 的系数正比于 $q$"这一结构不变。作者真正要用的是：$q\neq0$ 时 $J$ 的矩阵元在 $p'\to p$ 时非零，于是下文的自旋限制才能排除它。另注：原文此行 "TheP we will see" 显然是 "Then we will see" 的排印讹误。]

另一方面，利用洛伦兹不变性可以证明：对所有 $p'\neq p$，这些矩阵元在第一种情形下对大于 $1/2$ 的自旋为零，在第二种情形下对大于 $1$ 的自旋为零。
证明的出发点是先注意到：当 $p'\neq p$ 时，同时固定 $p$ 与 $p'$ 的洛伦兹子群是

% ===== 原文 p.37 =====
一个 $SO(2)$ 的拷贝（若我们在 $n$ 维时空中，则是 $SO(n-2)$；目前的考虑在低于四维时退化，因为那时 $SO(n-2)$ 是平凡的）。
人们只需证明：假定这些矩阵元非零，则 $\langle p',j|J^\nu|j,p\rangle$ 与 $\langle p',j|T^{\mu\nu}|j,p\rangle$ 的 $SO(2)$ 不变性蕴含自旋 $j$ 的绝对值 $\le 1/2$ 或 $\le 1$（分别对应两种情形）。
进行这一计算的一个方便办法，是进入这样的洛伦兹系（把时间坐标写在最前）：$p=(1,1,0,0)$，$p'=(1,-1,0,0)$。
同时固定 $p$ 与 $p'$ 的 $SO(2)$ 就是最后两个坐标的转动。
在该 $SO(2)$ 的生成元下，态 $|j,p\rangle$ 与 $|j,p'\rangle$ 的本征值分别为 $j$ 与 $-j$。
$|j,p'\rangle$ 态的这个负号至关重要，其来源是：它描述一个沿与 $|j,p\rangle$ 相反方向运动的粒子；相对于各自的运动方向，二者自旋相同，但若相对于某个固定轴来看，则自旋相反。
因此，$\langle p',j|J^\nu|j,p\rangle$ 或 $\langle p',j|T^{\mu\nu}|j,p\rangle$ 只有在如下情形下才可以非零：在 $SO(2)$ 对 $J$ 或 $T$ 的作用中，存在自旋（本征值）为 $-2j$ 的项。
由于 $J$ 的各分量在 $SO(2)$ 下按绝对值 $\le 1$ 的自旋变换，而 $T$ 的分量自旋 $\le 2$，我们便在两种情形下分别得到 $|j|\le 1/2$ 或 $|j|\le 1$，正如上面所断言。

% [Tutor Note: 关键是"同时固定 $p$ 与 $p'$ 的子群"取法：$p=(1,1,0,0)$ 与 $p'=(1,-1,0,0)$ 都沿 $x$ 方向（互为反向的光锥动量），故稳定化子恰为绕 $x$ 轴的转动 $SO(2)_{yz}$，其生成元是 $J_x$；对无质量态而言 $J_x$ 的本征值就是螺旋度（helicity）$j$，但对沿 $-x$ 方向运动的态，螺旋度按定义相对于自身动量方向，故 $J_x$ 的本征值为 $-j$。若误以为 $|j,p'\rangle$ 的本征值也是 $+j$，则"$-2j$"会算成"$0$"，结论（自旋上限）就完全消失。此外 $J^\nu$ 是矢量、在 $SO(2)_{yz}$ 下分解为自旋 $0,\pm1$ 三部分，$T^{\mu\nu}$ 是对称张量、分解为自旋 $0,\pm1,\pm2$，这正是上限 $1$ 与 $2$ 的来源。]

我们实际上可以把这个结果说得更精确一些。
迄今我们只用了表示论，但在量子场论中还有 CPT 定理，它在四维中蕴含：每个自旋为 $j$ 的无质量粒子都伴随一个自旋为 $-j$ 的无质量粒子。
因此自旋 $\pm1/2$ 会成对出现，同样自旋 $\pm1$ 也成对出现。

% [Tutor Note: 严格地说 CPT 是反幺正且反演动量的：它把 $|p,j\rangle$ 映到 $|p,-j\rangle$（差一个相位），因此"伴随"指的是同一动量、相反螺旋度的伙伴，而不是另一支独立的粒子。这也解释了为何下一段的 $n$ 维陈述中"旋量表示"是成对（$n$ 为偶数时有两个）出现的。]

在维数 $n>4$ 时，类似的推理给出如下结果。
无质量粒子的自旋由“小群”（little group）$SO(n-2)$ 的表示来分类。
若应力张量存在，则无质量粒子所允许的表示是旋量表示，\footnote{我们提醒：$n$ 为偶数时有两个这样的表示，$n$ 为奇数时只有一个。}以及基本 $(n-2)$ 维表示的外幂（包括平凡表示）。\footnote{当 $n-2$ 能被四整除时，中间那个外幂可分解为自对偶与反自对偶两部分，二者各自是实的；可以只有其一而无其二。}
除按平凡表示或旋量表示变换的无质量粒子外，整体荷都为零。
这一 $n$ 维的表述与四维中的描述（用 $j\in\mathbb{Z}/2$）有如下关系：$j=\pm1/2$ 对应两个旋量表示，而 $j=0$ 与 $j=\pm1$ 来自基本表示的外幂。

\subsection{红外极限}
% 原文 §2.5：Infrared limits

给定一个量子场论 $X$，我们想求解它，也就是想了解关于它最有趣的事情。
“求解”一个理论的一大部分工作，是确定它在红外流向什么。
相当常见的情形是，答案是“什么都没有”，即 $X$ 流向一个平凡理论。
这恰好发生在 $X$ 具有质量间隙（mass gap）之时，因为此时所有（欧氏）关联函数都指数衰减。
不过，证明一个给定的理论流向平凡理论，可能是一个相当深刻的结果。

非常常见的情形是，红外极限并不平凡，而是无质量粒子的自由理论，外加一个在红外趋于零的无关（irrelevant）相互作用。
事实上，我们将遇到的大多数最简单的例子都是这样。
注意，无关相互作用在紫外会被视为“不可重整化”；那些在紫外行为不良的微扰，恰是在流向红外极限时消失的微扰。

% ===== 原文 p.38 =====
当一个理论在红外是自由的，问题就变成：它是哪些无质量粒子的自由理论？
它们未必与原始拉格朗日量中的粒子相同。
（我们在 §1.8 中见过一个相对平凡的例子。）
事实上，正如我们在本讲前面所见，这个问题的答案可能依赖于我们所在的真空态。

红外极限平凡，是红外极限自由的一个特殊情形；此时物理谱中根本没有无质量粒子。

若一个理论中的耦合是通用的（generic），则无质量粒子必定有某个无质量的理由。
一个可能的理由是超对称，但现在我们不讨论它。
其它可能的理由如下。

自旋 $0$ 粒子在是 Goldstone 玻色子时无质量，即存在破缺的对称性时。
四维中自旋 $1/2$ 粒子在是手征费米子（chiral fermion）时无质量；它们的无质量性来自未破缺的手征对称性。
这简单地意味着：CPT 定理说，若有 $n$ 个自旋 $1/2$ 的无质量粒子，也就有 $n$ 个自旋 $-1/2$ 的此类粒子。
若存在未破缺的对称群 $G$，且自旋 $1/2$ 的无质量粒子按 $G$ 的表示 $R$ 变换，则由 CPT 定理，自旋 $-1/2$ 的无质量粒子按表示 $\bar R$（$R$ 的复共轭）变换。
若 $R$ 与 $\bar R$ 不同，则这个谱不能以 $G$ 不变的方式微扰，从而使费米子获得质量。

% [Tutor Note: 这一段是 't Hooft 式的"自然性"（naturalness）论证，作者一句话带过了其中的群论。注意"自旋 $1/2$"与"自旋 $-1/2$"在此指的是同一个 Weyl 费米子的两个螺旋度态（其 CPT 共轭），不是两组独立粒子；要给这个 Weyl 费米子质量，只能写 Majorana 型质量项，它要求 $R\otimes R$ 中含有 $G$ 的平凡表示，等价于 $R\cong\bar R$（$R$ 为实或赝实表示）。手征的定义正是 $R\not\cong\bar R$，故这样的项不存在，微扰无法给费米子质量。若误按 Dirac 质量项去要求 $R\otimes\bar R\supset\mathbf{1}$，则条件永远满足，整段论证就失效了——这是最容易出错的一步。另注：原文 "to a give masses" 是排印重复，应为 "to give masses"。]

上述理由实际上就是四维中在不用超对称、也不调节参数使粒子无质量的前提下，拥有自旋 $0$ 或 $1/2$ 无质量粒子的仅有的已知理由。
对自旋 $1$，情形有所不同。
在四（或更多）维中，无质量自旋 $1$ 粒子的庞加莱表示根本无法被微扰以使该粒子获得质量，除非存在一个无质量自旋 $0$ 粒子能通过 Higgs 机制与之结合，如本讲第 2 部分所述。
若我们的无质量自旋 $0$ 粒子是 Goldstone 玻色子，则平移它们的那些破缺对称性保证它们不能参与 Higgs 机制、从无质量谱中消失。
因此，只要无质量自旋 $0$ 粒子是 Goldstone 玻色子，且只要理论是红外自由的，无质量自旋 $1$ 粒子的存在仅凭庞加莱对称性就是稳定的。

这严重依赖于理论是红外自由的，因为我们把群论应用于单粒子态。
若相互作用在低能下仍然重要，不同粒子数的态可以“混合”，我们就不能只靠把庞加莱不变性应用于单粒子态而得出结论。

这类无质量粒子能有什么相互作用？
对自旋 $0$ 粒子（我们假定它们是 Goldstone 玻色子），没有相关（relevant）相互作用。
我们在本讲开头已经看到，仅有 Goldstone 玻色子时不存在相关或边缘（marginal）相互作用。
同样，在二维以上，仅有费米子时也不存在相关或边缘相互作用——我们在秋季学期探讨过这类问题——而用类似的论证并稍加小心，可以证明 Goldstone 玻色子与费米子之间不存在相关或边缘耦合。

自旋 $1$ 粒子性质不同，因为它们对应规范场，而规范场在红外\emph{可以}有相关相互作用。
若 $G$ 是规范群、$A$ 是联络，则拉格朗日量
\[
  \frac{1}{4g^2}\int |F_A|^2 d^4x
\]

% ===== 原文 p.39 =====
是非线性的。
它含有在四维中经典地为边缘（marginal）的耦合。
这些相互作用在量子力学上是相关还是无关，取决于 $\beta$ 函数的行为。
红外中无关的非线性将允许规范理论充当有效描述。
另一方面，若非线性在红外是相关的，则规范理论\emph{不是}答案。
（在 $\beta=0$ 的中间情形，我们会在红外得到一个非自由的理论。）

% [Tutor Note: 此处的“$\beta>0$”极易读反，务必与渐近自由的符号区分。按定义 $\beta(g)=\mu\,\partial g/\partial\mu$，$\beta>0$ 意味着耦合随能量升高而变大、在红外趋于零，这正是本段所需要的“红外自由”；而渐近自由（耦合在紫外趋于零）恰恰是 $\beta<0$。由于通常写 $\beta(g)=-b_0g^3/(16\pi^2)$，$\beta>0$ 等价于 $b_0<0$；而 $b_0=\frac{11}{3}C_2(G)-\frac{2}{3}\sum_R T(R)$（$T(R)$ 为表示 $R$ 的指标），故只有费米子的数目与表示足够大，才能把 $b_0$ 压到负值——这正是作者“足够多的费米子处于足够大的表示中”一句的定量含义。许多教科书写渐近自由时取 $b_0>0$，若照搬那套符号，就会把本段的判据说反。]

Goldstone 玻色子总是在规范对称性下不变的——一个作用于它们的规范群会破坏当初导致有 Goldstone 玻色子的那个整体对称性。
因此它们对 $\beta$ 函数没有贡献。
所以，若我们在红外看到无质量自旋 $1$ 场而没有为此调节参数，就应当预期：要么 $G$ 必须是阿贝尔的（abelian），要么必须有足够多的费米子处于 $G$ 的足够大的表示中，使得 $\beta>0$。
为解释它们的无质量性，这些费米子是手征的（即 $j=1/2$ 与 $j=-1/2$ 的态变换方式不同）——或者在规范群本身之下，或者在某个与规范群对易的未破缺整体对称群之下。

物理学中主要的未解决问题之一是：由于某个称为“规范等级问题”（gauge hierarchy problem）的东西，自然并不像上一段所暗示的那样通用。
这实际上是提示超对称可能在自然中重要的原因之一。

\subsection{参考文献与交叉引用}
% 原文 References and cross references

关于经典规范理论，见 [I-Classical Fields] 第 4 章。
规范理论的 BRST 量子化在本系列第 3 讲以及 [II-Strings] 第 4 讲中描述。
关于 Higgs 机制与红外极限的问题，见本系列课程末尾的习题（Exercises）、[I-Classical Fields] 第 2 章，以及 [I-Homework] 中超级作业（superhomework）的附录。
整体对称性破缺在第 1 讲中讨论。

关于弱耦合规范理论中 Higgs 机制的入门，见
\begin{quote}
S. Coleman, \textit{Aspects Of Symmetry}, Cambridge University Press, 1985.

S. Weinberg, \textit{The Quantum Theory Of Fields}, chapter 21, vol. 2, Cambridge University Press, 1996.
\end{quote}

% [Tutor Note: 文献代号 [I-Classical Fields]、[II-Strings]、[I-Homework] 指本丛书（\textit{Quantum Fields and Strings: A Course for Mathematicians}, AMS, 1999）第 I、II 卷中的对应章节；"superhomework" 是该书为数学背景读者设计的加长习题集，其 addendum 中补充了 Higgs 机制的红外讨论。引用时请保留原代号，不要译成中文。]
